\documentclass[12pt]{article}
\usepackage{amsmath,amssymb}
\usepackage[margin=1in]{geometry}

\begin{document}

\noindent
Professor E. Guazzelli

\vspace{1cm}

I am pleased to submit this manuscript on ``Linear Kelvin Wave Predictions in the $z\to 0$ Limit'' for your consideration as a Rapids article in the Journal of Fluids Mechanics. 

Linear wave theory is a timeless cornerstone of fluid mechanics, rich in theoretical insights, and still completely relevant for today's engineering and physics communities. However, the linear potential for a translating disturbance near the free surface is ill-posed, producing unbounded wave energy even far behind the disturbance. In this manuscript we propose a simple analytic solution to this problem by including the finite width of the disturbance directly into the formulation. The resulting regularized kernel produces physical uniformly finite predictions and clarifies fundamental wave generation mechanisms for finite-width geometries, providing insight into free-surface energy distribution and interference patterns. 

A second contribution of the manuscript is the development of a partitioned numerical method for evaluating the kernels. Extending recent work in oscillatory integral methods, the new approach is uniformly accurate, and provides a massive $10^4-10^5$ speed-up compared to adaptive quadrature. An open source Julia-language implementation is available and all of the results in the manuscript are fully reproducible. 

With both analytic and $\mu s$ numerical evaluations, the new theory enables practical use in high-speed design, control, and surrogate modeling, and establishes a robust linear foundation for studying Kelvin waves, making it an ideal fit for the Journal of Fluid Mechanics.

\vspace{1cm}

\noindent
Thank you for your time,\\ \\
Prof. Gabriel D Weymouth\\
TU Delft
\end{document}